% Proposed B.138. Owned source; central assembly and registers are root-owned.
\subsection*{The full topology and intersections of the actual unfilled edge groups}

\statement{B.138}{Proved}{The actual unfilled cyclic-edge subgroups have full induced topology, trivial whole completed intersection, and one exact ordinary translated overlap}
\label{stmt:B138}
Retain every marking of B.134--B.137. Put
\begin{equation}\label{eq:B138-data}
\begin{gathered}
V=R_L*_{K}R_R,\quad R_L=R^0,\quad R_R=p(R^0),\quad p|_K=\eta,\\
E_1=C\langle K,k\rangle C,\qquad \zeta=\eta(a),\\
B_i=\langle f,a^{-1}p(k)a\rangle,\qquad
B_j=\langle\eta(f),\zeta^{-1}k\zeta\rangle,\\
H_i=E_1*_{\langle f\rangle}B_i,\qquad
H_j=B_j*_{\langle\eta(f)\rangle}p(E_1).
\end{gathered}
\end{equation}
These are the actual embedded ordinary subamalgams. Both $H_i,H_j$
have their full induced profinite topology in $V$. Bars below mean
actual closures in the indicated full completions. In particular
\begin{equation}\label{eq:B138-whole-compact}
\overline{H_i}\cap\overline{H_j}=1
\qquad\text{inside }\widehat V.
\end{equation}
Their completed vertex and edge intersections are
\begin{equation}\label{eq:B138-compact-vertices}
\overline{H_i}\cap\widehat R_L=\widehat E_1,\quad
\overline{H_i}\cap\widehat R_R=\widehat B_i,\quad
\overline{H_i}\cap\widehat K=\overline{\langle f\rangle},
\end{equation}
with the corresponding actual $\theta$-images for $H_j$.
In the ordinary group there is also the specified whole translated
intersection
\begin{equation}\label{eq:B138-shifted}
H_i\cap\zeta H_j\zeta^{-1}=\langle k\rangle.
\end{equation}
The full two-family initial normal subgroup $N_*$ of B.134 meets
$\langle k\rangle$ trivially; its closed normal closure meets the
actual $\widehat{\mathbb Z}(k)$ trivially as well.

All these assertions precede the initial fillings. They do not assert
that $D_{\rm loc}$ survives in $V_*$, that its actual kernel is
$\kappa$-invariant, or that all translated intersections are controlled.
No completed analogue of \eqref{eq:B138-shifted} is used or asserted.

\begin{proof}
\emph{B.138.1. The complete marked tree input.}
Write, as in B.117,
\[
L=\langle K,k\rangle=H*_{T}P,\quad
H=\langle h,k\mid[h^2,k]=1\rangle,\quad T=\langle k\rangle,
\]
\[
P=\langle k\rangle\times F(a,d,e,V_4),\qquad
f=hk^{-1},\quad V_4=f^{-1}Bf.
\]
The actual color retraction gives $E_1\cap L=H$. The neck
retraction $r:L\to K$ fixes $K$, kills $k$, and maps $H$ onto
$\langle f\rangle\subset H$. Define
\begin{equation}\label{eq:B138-fork}
\begin{gathered}
\ell=afa^{-1},\quad q=\eta(\ell),\quad U=\langle k,q\rangle,\\
\gamma=afV_4f^{-1}a^{-1}f,\quad D=ded^{-1},\quad
w=\gamma Df^{-1}d^{-1},\\
q=wf^{-1}w^{-1},\qquad
\zeta=\gamma V_4^{-1}f^{-1}a^{-1},\qquad
\zeta B_j\zeta^{-1}=U.
\end{gathered}
\end{equation}
The last line consists of the literal whole-word identities checked
in B.137. Its unrestricted endpoint proof embeds the subdivided tree
of $\langle k\rangle*\langle q\rangle$ in the tree $X$ of $H*_TP$.
Call the image $Y$. Its fundamental segment joins $aH$ to $wH$.
The first core edge at $aH$ is $afT$, and its incident edges in $Y$
are exactly $k^nafT$. The $U$-stabilizer of $aH$ is $\langle k\rangle$;
every interior subdivision vertex has trivial stabilizer. These are
consequences of the equivariant embedded tree, not finite samples.

\emph{B.138.2. The specified translated intersection.}
The path from $H$ to $aH$ ends through $aT$. This is none of
$k^nafT$, since equality would give $k^nf\in T$, contrary to
the $h$-exponent. Thus the projection of $H$ to the convex tree
$Y$ is $aH$, and
\[
U\cap H=\langle k\rangle,\qquad E_1\cap U=\langle k\rangle.
\]
At the interior vertex $\gamma P$, the two core directions are
$1,D$ in $P/T$, whereas $V_4^{-1}$ differs from both in the free
leaf group. The vertex $\gamma V_4^{-1}H=\zeta aH$ lies on that
private branch and projects to $\gamma P$. Its $U$-stabilizer is
trivial. Consequently
\[
B_j\cap aHa^{-1}=1,\quad B_j\cap aE_1a^{-1}=1,
\quad B_i\cap\zeta p(E_1)\zeta^{-1}=1,
\]
where the last equality is the actual $\theta$-image.

For $u=e x\in U\cap E_1K$, with $e\in E_1,x\in K$, both $u,x$
belong to $L$, so $e\in H$. The retraction preserves $U,H$ and gives
\[
u r(u)^{-1}=e r(e)^{-1}\in U\cap H=\langle k\rangle.
\]
As $r(U)=\langle q\rangle$, this proves the exact set identity
\[
U\cap E_1K=\{k^nq^j:n,j\in\mathbb Z\}.
\]
Thus the common left incident edges for the invariant trees of
$H_i$ and $\zeta H_j\zeta^{-1}$ are precisely $k^nK$.
At the right endpoint apply $\theta$. If
$b\in B_j\cap aE_1a^{-1}K$, the same computation puts
$b r(b)^{-1}$ in $B_j\cap aHa^{-1}=1$. Hence $b\in K$,
so the only common right direction is the incoming edge.

The two ordinary invariant trees are embedded by B.137 and their
intersection is convex. Both contain the infinite star $\bigcup_nk^nK$;
all departures from it have just been excluded. Their intersection
is exactly that star. A common group element fixes its unique left
vertex and therefore lies in $E_1\cap U=\langle k\rangle$.
This proves \eqref{eq:B138-shifted}. The actual two regional base
windings give $V\to\mathbb Z^2$, kill $K$ and send $k$ to $(1,0)$.
All ten generators of $N_*$ are commutators. Winding therefore proves
the ordinary and completed cyclic survival assertions, without any
claim that the larger local quotients survive.

\emph{B.138.3. A cofinal all-modulus fork.}
For every $n\ge2$, define
\[
P_n=C_n(k)\times C_2(a)\times C_2(d)\times C_2(e)\times C_2(V_4),
\quad G_n=C_2(c_n)*P_n.
\]
There is a homomorphism of the entire $L$ to $G_n$ sending
$h\mapsto c_n$ and the leaf and $k$ generators to their indicated
images. Put $f_n=c_nk^{-1}$ and $H_n=\langle c_n,k\rangle=C_2*C_n$.
Then $G_n=H_n*_{C_n}P_n$.

Use subscript $n$ for the images of every whole word in
\eqref{eq:B138-fork}. The alternating leaf blocks of the fundamental
path become $a,V_4,a,e,d$, all nonidentity in $P_n/C_n$.
Their repetition cannot cancel through the intervening $f_n^{\pm1}$,
which remain outside $C_n$. Every nonzero power of $f_n$ is outside
$C_n$ by free-product normal form. At the first endpoint, for every
nonzero residue $j\pmod n$,
\[
f_n^{-1}k^jf_n=kc_nk^jc_nk^{-1}\notin C_n.
\]
The endpoint coset injections therefore embed the subdivided tree
of $C_n(k)*\mathbb Z(q)$, proving
\begin{equation}\label{eq:B138-Un}
U_n=\langle k_n,q_n\rangle\cong C_n*\mathbb Z,\qquad
\ker(U\to U_n)=\langle\!\langle k^n\rangle\!\rangle_U.
\end{equation}
At the interior fork the three directions remain $1,e,V_4$ and are
distinct. The entire remaining tail to $\zeta_nH_n$ is reduced.
Projection to the interior vertex gives, for every $n\ge2$,
\begin{equation}\label{eq:B138-finite-target-intersection}
U_n\cap\zeta_nH_n\zeta_n^{-1}=1.
\end{equation}

\emph{B.138.4. The precise finite-factor topological input.}
A free product or amalgam of finite groups is virtually free here
by a direct finite permutation construction. Take enough copies of
the regular vertex actions so that their restrictions to the common
finite edge have the same number of free orbits, and identify those
orbits equivariantly. The resulting finite quotient is injective on
both vertices. Its kernel acts freely on the ordinary tree and has
finite quotient graph, hence is a finite-rank free group.

Finitely generated subgroups of a finite-rank free group are closed:
their generating loops lie in a finite core of the Schreier cover.
Adjoin the finite path of any word outside the subgroup, then complete
the finite labeled graph to a finite cover by pairing the missing
incoming and outgoing slots of each label. No old vertex is identified,
so that word remains excluded. Finite-index cores transfer this to
virtually free groups. Every finite-index subgroup of a finitely
generated subgroup is also finitely generated and closed; separating
its finitely many other cosets and taking an ambient normal core
shows that every finitely generated subgroup has its full induced
profinite topology.

We additionally use Ribes--Zalesskii's closed-product theorem for
free products \cite[Theorem~3.13]{BRibesZalesskii2004FreeProducts},
with the class of all finite groups. Finite factors satisfy its
hypotheses, so products of two finitely generated subgroups of $G_n$
are closed. Their Lemma~3.1 supplies the corresponding finite-index
passage. The closed-subgroup hypotheses of that theorem are met by
the preceding explicit argument.

Here is the exact intersection consequence. For finitely generated
$A,B\le G_n$, take $x\in\overline A\cap\overline B$ and a finite
quotient $\pi:G_n\to F$. Choose $a\in A,b\in B$ with image
$\widehat\pi(x)$, and put $A_0=A\cap\ker\pi$, $B_0=B\cap\ker\pi$.
Then $x\in a\overline{A_0}\cap b\overline{B_0}$, whence
$b^{-1}a\in\overline{B_0A_0}$. Closedness of this product gives
$b^{-1}a=b_0a_0$ with $b_0\in B_0,a_0\in A_0$. The element
$bb_0=aa_0^{-1}\in A\cap B$ has the prescribed finite image.
As $\pi$ was arbitrary,
\begin{equation}\label{eq:B138-completion-intersection}
\overline A\cap\overline B=\overline{A\cap B}.
\end{equation}
This proves the needed completion comparison; it is not assumed.

\emph{B.138.5. Actual compact vertex comparisons.}
Given any finite target of the actual free pair $U$, choose $n\ge2$
divisible by the order of the image of $k$. By \eqref{eq:B138-Un},
it factors through $U_n$. Full topology of $U_n$ in $G_n$ supplies
a finite $G_n$ quotient refining that target. Consequently $U$ has
full topology in $L$, and the maps $\widehat U\to\widehat G_n$
jointly separate every completed element. Equations
\eqref{eq:B138-finite-target-intersection}--\eqref{eq:B138-completion-intersection}
therefore give $\overline U\cap\zeta\widehat H\zeta^{-1}=1$.
Conjugating back gives $\overline{B_j}\cap\widehat H=1$ and
full topology of $B_j$ in $R_L$.

The continuous color retraction $\widehat R_L\to\widehat L$
fixes $\widehat L$ and sends $\widehat E_1$ onto $\widehat H$.
It follows that
\begin{equation}\label{eq:B138-regional-compact}
\widehat E_1\cap\overline{B_j}=1,\qquad
\widehat E_1\cap\widehat L=\widehat H.
\end{equation}
If $b=e x\in\overline{B_j}\cap(\widehat E_1\widehat K)$,
compactness supplies the indicated factors; as $b,x\in\widehat L$,
$e\in\widehat H$. The continuous neck retraction preserves both
$\overline{B_j}$ and $\widehat H$, so
$b r(b)^{-1}=e r(e)^{-1}=1$. Thus
\begin{equation}\label{eq:B138-product-comparison}
\overline{B_j}\cap(\widehat E_1\widehat K)
=\overline{B_j}\cap\widehat K
=\overline{\langle\eta(f)\rangle}.
\end{equation}
The actual $\theta$ supplies the right-vertex statements with its
entire neck marking preserved.

\emph{B.138.6. Extend an arbitrary finite target of the whole edge group.}
Let $\psi:H_i\to F_0$ be any finite target. Full vertex topology,
proved above for $B_i$ and by the color retraction for $E_1$, supplies
finite quotients $q_L,q_R$ of $R_L,R_R$ refining the two restrictions
of $\psi$. The actual identified-neck retractions satisfy
\[
r_L(E_1)=\langle f\rangle\subset E_1,\qquad
r_R(B_i)=\langle f\rangle\subset B_i.
\]
Choose $M\triangleleft_{\rm fi}K$ below both restrictions of
$\ker q_L,\ker q_R$, and augment each quotient by its retraction
modulo $M$. Their finite images $A,B$ have exactly the same edge
$C=K/M$. This defines $V\to G=A*_CB$.

Let $A_1,B_1$ be the respective images of $E_1,B_i$ and $C_1$
the image of $\langle f\rangle$. Then
\[
A_1\cap C=C_1=B_1\cap C.
\]
For example, if the image of $e\in E_1$ equals that of $x\in K$,
the extra coordinate gives $xM=r_L(e)M=f^nM$. Since $M$ is below
the original edge kernel, the entire image equals that of $f^n$.
The converse and right equality follow identically. Ordinary normal
form embeds $J=A_1*_{C_1}B_1$ in $G$, and the two restrictions of
$\psi$ factor through these finite vertices and agree on $C_1$.
Hence $\psi$ factors through the whole $J$, including its mixed words.
The finitely generated $J$ has full topology in the virtually free $G$.
A finite $G$ quotient refining this target thus gives a finite $V$
quotient refining $\psi$. This proves full topology of $H_i$, and
$\theta$ proves it for $H_j$.

The same arbitrary-target construction proves
\eqref{eq:B138-compact-vertices}. For instance, an element of
$\overline{H_i}\cap\widehat R_L$ maps in $\widehat G$ into
$\overline J\cap A$. Since $A$ is finite and $J$ is closed in $G$,
this intersection is $J\cap A=A_1$. Its image in every prescribed
finite $H_i$ target therefore lies in the image of $E_1$, forcing
membership in $\widehat E_1$. The right and edge versions use $B,C$.
These equalities also make the vertex- and edge-coset maps from the
standard completed $H_i$ tree into the completed $V$ tree injective;
compactness makes them embeddings. No finite-geodesic argument on
that profinite tree will be used below.

\emph{B.138.7. Two profinite coproduct lemmas.}
First, for profinite groups with a common continuous retract,
\begin{equation}\label{eq:B138-retract-coproduct}
(N_A\rtimes K)\amalg_K(N_B\rtimes K)
       \cong(N_A\amalg N_B)\rtimes K.
\end{equation}
The continuity of the right action merits proof. Each $k\in K$
extends its two factor actions to the coproduct $N$. For a finite
continuous map $\varphi:N\to F$, joint continuity on a compact
factor gives a neighborhood of $1$ in $K$ fixing the entire restricted
finite map: cover that factor by finitely many continuity neighborhoods.
An open subgroup fixing both restricted maps fixes $\varphi$ on $N$.
Its orbit under $K$ is finite. Intersect those finitely many kernels
to obtain a $K$-invariant open normal subgroup below $\ker\varphi$.
The induced finite action is continuous by the same factor argument.
These subgroups are cofinal, proving joint continuity. Both inverse
maps in \eqref{eq:B138-retract-coproduct} now follow from the
universal properties and fix the kernel factors and $K$.

Second, for arbitrary closed subgroups $A_1,A_2\le A$ and
$B_1,B_2\le B$ in a free profinite coproduct,
\begin{equation}\label{eq:B138-factorwise}
\overline{\langle A_1,B_1\rangle}\cap
\overline{\langle A_2,B_2\rangle}
=\overline{\langle A_1\cap A_2,B_1\cap B_2\rangle}.
\end{equation}
Fix an arbitrary finite target $q:A\amalg B\to D$.
There is a finite quotient $A\to A_0$ refining $q|_A$ such that
the intersection of the two subgroup images projects into
$q(A_1\cap A_2)$. Otherwise, in the compact space $A_1\times A_2$,
the pairs equal modulo each finer open normal kernel but with common
$q$-value outside this set form nonempty closed sets with the
finite-intersection property. A limiting pair is equal in $A$,
a contradiction. Do the same in $B$.

In the ordinary finite-factor free product $G_0=A_0*B_0$, the
two image subgroups are $H_0=A_{10}*B_{10}$ and $J_0=A_{20}*B_{20}$.
They are finitely generated and closed, and ordinary normal form gives
\[
H_0\cap J_0=(A_{10}\cap A_{20})*(B_{10}\cap B_{20}).
\]
The closed-product consequence \eqref{eq:B138-completion-intersection}
applies in $\widehat G_0$. Thus every common compact element has
$q$-image in the subgroup generated by $q(A_1\cap A_2)$ and
$q(B_1\cap B_2)$. All finite $q$ prove \eqref{eq:B138-factorwise}.
No finite-generation assumption on the original compact factors or
subgroups is present.

\emph{B.138.8. The whole completed intersection.}
The actual vertex retractions agree on $K$, giving
$r_V:\widehat V\to\widehat K$. Its images on the two edge groups
are the cyclic closures of $f$ and $\eta(f)$. They meet trivially.
Indeed the literal word is
\[
\eta(f)=P f^{-1}P^{-1},\qquad
P=aBfde d^{-1}f^{-1}d^{-1}.
\]
For every $n\ge2$, kill $B,d,e$ and map $K$ to $C_n(f)*C_2(a)$.
Then $P\mapsto a$. The two relevant finite cyclic subgroups are
disjoint by ordinary free-product normal form and remain so in the
full completion of this residually finite group. A common $f^\lambda$
therefore has $\lambda=0\pmod n$ for every $n$, so is one.

Put $N_L=\ker(\widehat R_L\to\widehat K)$ and similarly $N_R$.
Equation \eqref{eq:B138-retract-coproduct} gives
$\ker r_V=N_L\amalg N_R$. The full topology proved in B.138.6
identifies each completed edge group with its actual proper cyclic
amalgam. Applying the same retract lemma to it identifies its kernel
as the closed subgroup generated by its two vertex kernels. These are
\[
\begin{array}{c|cc}
 &N_L&N_R\\ \hline
\ker(r_V|_{\overline{H_i}})&
 \ker(r_L|_{\widehat E_1})&\ker(r_R|_{\widehat B_i})\\
\ker(r_V|_{\overline{H_j}})&
 \ker(r_L|_{\widehat B_j})&\ker(r_R|_{\widehat{p(E_1)}}).
\end{array}
\]
The two kernels in either column intersect trivially by
\eqref{eq:B138-regional-compact} and its mirror. The factorwise lemma
\eqref{eq:B138-factorwise} therefore gives trivial intersection of
the two edge kernels. Any element of $\overline{H_i}\cap\overline{H_j}$
already lies in both kernels by the cyclic comparison. This proves
\eqref{eq:B138-whole-compact}.

The arguments do not impose either initial relator family. A trivial
unshifted intersection does not identify the group generated by the
two edge subgroups, classify their translated overlaps, or identify
the restriction of their ambient normal closures. Both fillings and
the ordinary and compact matching criteria of B.134 remain necessary;
no closure/intersection or ordinary/compact kernel equality is added.
\end{proof}

\begingroup\raggedright\footnotesize
The complete frozen proof records are
\path{work_DE/next_marginal_source/ambient_shifted_overlap28.md},
\path{work_DE/next_marginal_source/compact_vertex_overlap28.md},
\path{work_DE/next_marginal_source/compact_actual_edge_embedding28.md},
and \path{work_DE/next_marginal_source/compact_whole_edge_intersection28.md}.
All root/A/CE complete reviews and the whole-input B.137 fork checker
are retained. The complete primary Theorem3.13/proof and Lemma3.1
proof, author publication metadata, and the failed publisher retrieval
are recorded in
\path{work_DE/next_marginal_source/RibesZalesskii2004_primary_audit28.md}.
No finite quotient, period or translated-word census replaces any
all-target or all-exponent argument above.
\par\endgroup
